\documentclass[12pt]{article}
\usepackage[utf8]{inputenc}
\usepackage[T1]{fontenc}
\usepackage{hyperref}
\title{Warji Language Resources for K-12 Curriculum Development}
\author{Compiled by Claude Sonnet 5.5 (claude-sonnet-5-5) for Andr\'as Kornai}
\date{}
\begin{document}
\maketitle

\section{Introduction}
Warji (ISO 639-3 wji; Glottocode warj1253), also called Sirzakwai or Warjawa, is a West Chadic (Afro-Asiatic) language of Bauchi State, Nigeria. It belongs to the North Bauchi group. Speaker numbers are given as roughly 75,000--78,000. Speakers are reported to be shifting toward Hausa. Documentation is sparse.

\section{Primary Educational Resources}
No primers, textbooks or school programmes in Warji were found in search results. The language is not known to be taught in schools.

\section{Linguistic and Reference Materials}
\begin{itemize}
\item Overview, including the elaborate noun plural system: \url{https://en.wikipedia.org/wiki/Warji_language}
\item ASJP wordlist entry for Warji: \url{https://asjp.clld.org/languages/WARJI}
\item Max Planck Institute numerals page for Warji: \url{https://lingweb.eva.mpg.de/channumerals/Warji.htm}
\item Endangered Languages Project entry and bibliography: \url{https://endangeredlanguages.com/node/96619}
\item Roger Blench's work-in-progress Warji (Sirzakwai) wordlist with Hausa and English glosses is reported on Academia.edu as ``Warji Dictionary'' (academia.edu, item 3764867). It appeared only in a search summary, not as a listed result link, so it is not in the JSON.
\item Bibliographic leads not checked directly: Skinner (1977) North Bauchi survey; Jungraithmayr (1965/66), \emph{Specimens of the Pa'a and Warja Languages}.
\end{itemize}

\section{Religious and Community Resources}
No scripture translation or community organisation for Warji was found in these searches.

\section{Digital and Audio Resources}
No audio, video, corpus or software resources were found beyond the wordlist databases above.

\section{Curriculum Development Context}
There is no full grammar, no downloadable dictionary and no literacy material that could be confirmed. A curriculum would need fresh documentation work with the community, using the Blench and Skinner wordlists as starting points and Hausa as the likely language of schooling.

\end{document}